\chapter{Atualização OTA, publicação e retorno de versão}\label{cap:ota}

\section{Objetivo e limites da atualização}
OTA significa \emph{Over-the-Air}: substituir o firmware sem repetir a gravação por cabo. No NEXUS, apenas o mestre obtém as imagens da Internet. Ele instala sua própria imagem por HTTPS ou transmite pela UART a imagem destinada a um satélite. A decisão de \emph{instalar} continua condicionada a uma autorização do operador. Uma versão encontrada, uma ordem aceita pelo servidor web, um bloco UART confirmado e uma imagem gravada são acontecimentos distintos; nenhum deles comprova, isoladamente, que o novo programa iniciou e permaneceu funcional.

Há três firmwares independentes: mestre ESP32-S3 N16R8, satélite da Equipe e satélite de Visitantes. Cada um usa duas partições de aplicação, chamadas \texttt{ota\_0} e \texttt{ota\_1}. A atualização grava a partição que não está executando e conserva a imagem anterior até a troca de boot. Essa escolha protege a operação contra download interrompido e permite ao \emph{bootloader}, o pequeno programa que escolhe o aplicativo na partida, retornar à imagem anterior se a nova falhar no primeiro ensaio. A instalação inicial \textbf{precisa ser feita por cabo} para gravar também a tabela de partições; simplesmente atualizar um firmware antigo que use uma única partição não cria o espaço OTA.

A tabela~\ref{tab:particoes_nexus} apresenta os deslocamentos e tamanhos exatos da flash do mestre e dos satélites. Cada slot de aplicação tem 6\,MiB no mestre e 1\,920\,KiB nos satélites; \texttt{otadata} mantém a seleção de partida, enquanto \texttt{nvs} conserva preferências fora das imagens. A imagem nova deve caber \emph{inteira} no slot inativo. O mestre verifica o limite antes de baixar e o satélite rejeita no comando de início um tamanho maior que o seu slot. Não se deve gravar um binário ou tabela preparados para o N16R8 em um SuperMini de 4\,MiB.

Para instalar essa tabela pela primeira vez, conecte cada placa por USB, escolha a porta correspondente no PlatformIO e grave os ambientes corretos. As instalações posteriores podem seguir por OTA, desde que as partições já estejam presentes:
\begin{lstlisting}[caption={Primeira gravação com a tabela de partições OTA}]
pio run -d firmware/mestre-s3 -e mestre-s3-n16r8 -t upload
pio run -d firmware/equipe-s3 -e esp32-s3-devkitm-1 -t upload
pio run -d firmware/visitantes-s3 -e esp32-s3-devkitm-1 -t upload
\end{lstlisting}

\section{Versões, manifesto e conjunto publicável}
A versão de cada aplicativo tem forma \texttt{MAJOR.MINOR.PATCH}. Incrementar \emph{major} anuncia incompatibilidade deliberada que requer migração ou atualização coordenada; \emph{minor} acrescenta função preservando a compatibilidade prevista; \emph{patch} corrige comportamento sem mudar o contrato. O número na etiqueta do GitHub Release identifica o \emph{conjunto} publicado. As três versões internas podem ser diferentes: uma publicação nova pode transportar um satélite alterado e dois binários inalterados. Comparar apenas a etiqueta do release produziria atualizações desnecessárias; por isso, o mestre compara separadamente a versão informada no manifesto com a versão em execução de cada alvo.

O arquivo \texttt{manifesto-firmwares.json} é o índice desse conjunto. Seu formato atual é \texttt{versao\_formato: 2}; contém \texttt{firmwares.mestre-s3}, \texttt{firmwares.equipe-s3} e \texttt{firmwares.visitantes-s3}. Para cada alvo há versão, tamanho em bytes, endereço HTTPS do binário e SHA-256. SHA-256 é um resumo criptográfico de 256 bits: uma mudança de conteúdo altera o resumo com probabilidade extremamente alta. O firmware exige endereço \texttt{https://}, resumo hexadecimal válido, versões interpretáveis e os três alvos presentes; limita o manifesto recebido a 8\,192 bytes. A versão dos satélites também precisa caber no campo de versão do protocolo UART. A etiqueta e os endereços devem corresponder ao mesmo conjunto de anexos. \textbf{HTTPS autentica o servidor por certificado; SHA-256 compara o arquivo recebido com o valor anunciado pelo manifesto.} O manifesto público, por si só, não é uma assinatura digital independente da conta e do canal de publicação; proteger permissões do repositório e do release continua necessário.

O mestre consulta \path{releases/latest/download/manifesto-firmwares.json} do repositório configurado. O release destinado à frota deve, portanto, ser o mais recente atendido por esse caminho e conter os quatro anexos completos: três \texttt{.bin} e o manifesto. O firmware atual não fornece credenciais para releases privados. Se a publicação for privada, incompleta ou apontar a binário ausente, a consulta ou a transferência falha sem instalar uma imagem parcial. O cabeçalho \texttt{Date} de uma resposta HTTPS válida pode acertar o relógio UTC dos logs no SD quando falta uma hora confiável; isso não substitui a validação TLS.

\subsection{Preparação local e publicação}
O estado de referência é \texttt{versoes/estado\_versoes.json}. Ele registra a última etiqueta, as três versões e impressões das fontes. O preparador considera código compartilhado em todos os alvos afetados e também HTML, CSS, JavaScript, geradores e páginas C geradas dos satélites. Arquivos temporários de compilação não devem determinar a versão. Depois de escolher conscientemente \texttt{patch}, \texttt{minor} ou \texttt{major}, execute na raiz do repositório:
\begin{lstlisting}[caption={Preparação local do conjunto OTA}]
python ferramentas/preparar_release_ota.py --bump minor
\end{lstlisting}
Sem \texttt{--bump}, o script oferece um menu; \texttt{--forcar equipe-s3} ou outro alvo só deve ser usado quando é necessário reconstruir e publicar nova versão sem mudança detectada. A preparação regenera painéis, detecta fontes alteradas, compila os alvos necessários e confere a versão \emph{embutida} em cada \texttt{.bin}. Imagens anteriores só são reaproveitadas quando correspondem ao estado registrado. Confirme os quatro arquivos em \texttt{versoes/saida/}, as versões internas, o tamanho e o SHA-256 de cada binário, bem como a etiqueta única indicada no manifesto. Não edite o número de versão somente no JSON: isso criaria divergência com a descrição do aplicativo gravada no binário.

Preparar arquivos locais não os publica. O caminho normal de publicação é o workflow \path{.github/workflows/publicar-firmware-mestre.yml}. Faça commit do código, do estado de versões e das páginas web geradas, envie esse commit, crie a etiqueta \texttt{vX.Y.Z} \emph{nele} e envie a etiqueta. O workflow recompila, rejeita páginas geradas desatualizadas e publica os quatro anexos no GitHub Release. Confira o resultado em Actions e Releases, inclusive as notas de migração quando houver incompatibilidade. Uma tag criada antes do commit correto, um workflow falho ou um release com anexos misturados não constitui uma versão OTA utilizável. Não reutilize uma etiqueta para conteúdo diferente, pois isso torna ambíguo o diagnóstico de qual imagem foi instalada.

\section{Consulta e autorização no equipamento}
O intervalo padrão de consulta é de 15 minutos; a configuração admite de 1 minuto a 7 dias. O mestre abre uma sessão Wi-Fi de finalidade OTA, busca o manifesto por HTTPS, compara as versões e fecha a sessão quando possível. O prazo HTTP configurado é de 20 segundos por operação. O Wi-Fi não fica ligado continuamente apenas para procurar versões. Uma sessão de outra finalidade em andamento impede uma verificação concorrente, evitando disputa pelo rádio, pela conexão e pelo estado da operação.

O satélite somente aparece como candidato se estiver respondendo recentemente, anunciar capacidade OTA UART e tiver versão inferior à do manifesto. O mestre considera a resposta de telemetria recente até 5 segundos para avaliar o enlace. A ausência de resposta, capacidade ou versão nova é mostrada como indisponibilidade, falta de suporte ou atualização desnecessária, conforme o caso. Uma oferta de atualização permanece aguardando autorização por \textbf{5 minutos}; expirado o prazo, é descartada sem alterar o firmware em execução. Uma nova consulta pode oferecer a mesma versão novamente. Esse prazo evita que uma autorização tardia instale uma decisão tomada com informação antiga sobre enlace, versão ou publicação.

No terminal serial do mestre a 115\,200 baud, a opção \texttt{2} abre o submenu OTA: \texttt{2} verifica agora; \texttt{3}, \texttt{4} e \texttt{5} autorizam, respectivamente, mestre, Equipe e Visitantes; \texttt{6} atualiza todos; \texttt{7} cancela; \texttt{8} altera o intervalo de consulta; \texttt{9} altera o intervalo de telemetria durante OTA. A página \emph{Comandos} do ponto de acesso da Equipe oferece autorizações equivalentes. Quem se associa ao AP Equipe pode enviar a ação HTTP sem uma chave administrativa adicional. A ponte Equipe--mestre devolve resposta do comando, mas o andamento deve ser consultado no estado OTA. Uma resposta de pedido aceito significa que o serviço registrou a autorização; \textbf{não significa que o download começou, que os bytes chegaram ou que o novo aplicativo foi confirmado}. A ação para \emph{todos} depende dos três alvos; se um satélite estiver ausente, pode-se selecionar apenas o alvo que responde, mantendo a verificação de versão e do estado final.

Na opção \emph{Atualizar todos}, a ordem é Equipe, Visitantes e mestre. Alvos já atualizados são pulados. Falha de um satélite interrompe a sequência antes de substituir o mestre. O motivo é operacional: o mestre antigo permanece disponível para diagnóstico e para repetir a tentativa por alvo. Uma atualização de só um alvo é adequada quando a mudança é compatível; para alteração \emph{major} do protocolo, as notas do release devem especificar a ordem segura de migração e a compatibilidade de versões intermediárias.

\section{Instalação do mestre por HTTPS}
Após a autorização, o mestre escolhe com \texttt{esp\_ota\_get\_next\_update\_partition} o slot inativo. \texttt{esp\_https\_ota} baixa e grava a imagem nesse slot, usando validação do certificado TLS pelo conjunto de raízes configurado. Antes de qualquer troca de boot, compara a versão declarada \emph{dentro} do binário com a versão prometida no manifesto, confere o tamanho informado quando disponível, rejeita imagem que exceda o slot, confere a quantidade total efetivamente recebida e calcula o SHA-256 da partição gravada. O gerenciador só chama \texttt{esp\_https\_ota\_finish} após esses testes; nessa etapa a partição passa a ser a escolhida para a próxima partida. Em caso de erro ou cancelamento antes disso, aborta a operação e a partição anterior continua executando.

O serviço registra bytes recebidos, porcentagem em décimos, taxa e estimativa de tempo restante. Quando a gravação termina com sucesso, informa estado \emph{aplicado} e reinicia após aproximadamente 3 segundos. Esse estado indica \textbf{imagem escrita, íntegra e selecionada para a primeira partida}; ainda não demonstra que as tarefas, Wi-Fi, UART e demais serviços funcionam com o novo firmware. A confirmação funcional ocorre depois do reboot e da janela de observação descrita adiante.

\section{Instalação dos satélites: HTTPS até o mestre, UART até o destino}
Os satélites não precisam acessar o GitHub. O mestre baixa por HTTPS o binário específico do alvo e o encaminha em blocos pela UART dedicada. O serviço usa três buffers reutilizáveis de 4 KiB, preferencialmente na PSRAM e, se ela não estiver disponível, tenta alocá-los na RAM interna. Uma tarefa pode preencher o próximo buffer pela rede enquanto outra transmite e aguarda confirmação do bloco anterior. Isso evita manter o firmware inteiro na RAM ou no cartão SD. O quadro entregue ao driver UART precisa de memória interna apropriada ao driver; a PSRAM serve como área de encadeamento, não como substituição irrestrita de buffers do periférico.

O mestre escolhe 4\,096 bytes por bloco quando o satélite anuncia essa capacidade. Para firmware compatível que só anuncia OTA UART, usa 2\,048 bytes. Durante o envio, a telemetria não é simplesmente desligada: o intervalo padrão passa a 2 segundos e pode ser ajustado entre 0,5 e 10 segundos. Isso conserva observabilidade básica, mas deixa mais tempo de UART para a imagem. A velocidade real depende de rede, UART, gravação da flash e confirmações; o percentual é calculado a partir dos \emph{bytes confirmados} pelo satélite, não apenas baixados pelo mestre.

\begin{table}[htbp]\centering\small
\caption{Quadro OTA UART versão 1 dos satélites}\label{tab:quadro-ota-satelite}
\begin{tabularx}{\textwidth}{@{}l l X@{}}\toprule
Campo & Tamanho & Finalidade\\\midrule
Marcador \texttt{A5 5A} e versão & 3 bytes & Identificam a família de quadro e o formato.\\
Tipo & 1 byte & Iniciar, dados, finalizar, cancelar ou resposta.\\
Sessão & 4 bytes & Distingue uma transferência de outra.\\
Sequência & 4 bytes & Ordena quadros e permite reconhecer repetição.\\
Deslocamento & 4 bytes & Início do bloco, ou próximo byte esperado na resposta.\\
Tamanho da carga & 2 bytes & Permite delimitar até 4\,096 bytes.\\
Carga & variável & Dados da imagem ou metadados e resultado.\\
CRC-32 & 4 bytes & Detecta alteração acidental do quadro recebido.\\\bottomrule
\end{tabularx}\end{table}

O cabeçalho compactado soma 18 bytes; campos multibyte são transmitidos em \emph{little-endian}. O CRC-32 verifica cada quadro, enquanto SHA-256 verifica o arquivo completo: o primeiro ajuda a decidir se um bloco deve ser repetido; o segundo impede selecionar uma imagem final diferente da anunciada. A carga \emph{INICIAR} tem 60 bytes: tamanho total de 32 bits, SHA-256 binário de 32 bytes e versão esperada em 24 bytes. A resposta contém código, fase e erro ESP-IDF; o cabeçalho devolve sessão, sequência e próximo deslocamento. As fases incluem ocioso, recebendo, validando, pronto para reiniciar, cancelado e falha. O identificador de sessão é sorteado pelo mestre; um cancelamento antigo não pode encerrar uma nova transferência.

No início, o satélite verifica formato, versão, tamanho e espaço do slot inativo, abre a gravação por \texttt{esp\_ota\_begin} e responde que está recebendo. Em cada bloco, exige sessão, sequência e deslocamento esperados antes de chamar \texttt{esp\_ota\_write}. Se a resposta se perder e o mestre repetir o último bloco já gravado, o satélite reconhece a repetição e confirma o mesmo próximo deslocamento sem gravar novamente. Sequência ou posição incoerente recebe erro explícito. O mestre valida sessão, sequência, fase e deslocamento da resposta: uma confirmação de outro bloco não avança o progresso.

Cada envio admite até quatro tentativas. O prazo de resposta é 1,5 segundo para dados, 10 segundos para iniciar e 15 segundos para finalizar. Esses prazos maiores envolvem apagar e preparar a partição ou completar a verificação final. Sem dados da sessão por 15 segundos, o satélite aborta a transferência parcial por inatividade; a verificação continua mesmo se trafegarem outros tipos de quadro na UART. Se o pipeline interno do mestre ultrapassa seu próprio prazo, o serviço o bloqueia até uma reinicialização segura em vez de arriscar reutilizar buffers ou respostas de uma operação indeterminada.

O comando \emph{FINALIZAR} só é aceito quando todos os bytes esperados foram gravados. O satélite chama \texttt{esp\_ota\_end}, calcula o SHA-256 da partição, compara com o resumo anunciado, lê a descrição do aplicativo e compara sua versão com a esperada. Só então seleciona a partição de boot, transmite a fase \emph{pronto para reiniciar} e reinicia cerca de 1 segundo depois. Ele tenta emitir essa resposta três vezes; se nenhuma transmissão for bem-sucedida, tenta restaurar a partição anterior e cancela o reinício. Se o mestre não recebe confirmação coerente, considera a tentativa falha, mesmo que alguns bytes tenham sido recebidos. Depois do reinício, espera por até 15 segundos até que uma resposta de telemetria recente anuncie a versão esperada. Esse é o critério de \emph{conclusão da transferência} no mestre; a confirmação interna de estabilidade do satélite ainda ocorrerá no prazo de 2 minutos.

\section{Estados da imagem e duas janelas de proteção}
O \texttt{otadata} registra qual slot deve iniciar e o estado da imagem. No mecanismo de rollback do ESP-IDF, a imagem recém-selecionada entra como \texttt{NEW}; na primeira partida passa a \texttt{PENDING\_VERIFY}. Se o aplicativo reiniciar antes de confirmar, o bootloader marca essa tentativa como abortada e procura a imagem válida anterior. Se os serviços essenciais falharem durante a condição pendente, o aplicativo solicita \path{esp_ota_mark_app_invalid_rollback_and_reboot}. Após iniciar e funcionar por 2 minutos, chama \path{esp_ota_mark_app_valid_cancel_rollback}, que marca a imagem válida. O \texttt{otadata} tem 8 KiB, organizados pelo ESP-IDF para conservar os dados de seleção entre gravações e reinícios; o binário continua nos slots \texttt{ota\_0}/\texttt{ota\_1}.

\begin{figure}[htbp]\centering
\begin{tikzpicture}[node distance=7mm, >={Latex},
estado/.style={draw=nexus,fill=pale,rounded corners,align=center,text width=9.6cm,minimum height=9mm,font=\small}]
\node[estado] (anterior) {Imagem anterior válida; slot inativo disponível};
\node[estado,below=of anterior] (escrita) {Download e gravação no slot inativo; tamanho, versão e SHA-256 conferidos};
\node[estado,below=of escrita] (novo) {Slot novo selecionado; bootloader inicia imagem em \texttt{PENDING\_VERIFY}};
\node[estado,below=of novo] (dois) {Serviços essenciais ativos e 2 minutos de observação};
\node[estado,below=of dois] (cinco) {Imagem \texttt{VALID}; guarda RTC acompanha reinícios anormais até 5 minutos estáveis};
\node[estado,below=of cinco] (fim) {Observação encerrada; imagem em operação normal};
\draw[->,thick] (anterior)--(escrita);
\draw[->,thick] (escrita)--(novo);
\draw[->,thick] (novo)--(dois);
\draw[->,thick] (dois)--(cinco);
\draw[->,thick] (cinco)--(fim);
\end{tikzpicture}
\caption{Caminho normal de uma imagem OTA. Falhas, cancelamento e retorno à imagem anterior estão descritos na tabela~\ref{tab:falhas-ota}.}\label{fig:estados-ota}
\end{figure}

O mestre e os dois satélites registram a partida antes de completar os serviços. Durante o estado pendente, falha considerada essencial no início dos serviços provoca rejeição da imagem e pedido de retorno; em especial, falhas de Wi-Fi ou do serviço OTA comprometem a capacidade de recuperação. Depois de inicializar os serviços, a aplicação aguarda 120\,000 ms, confirma a imagem em execução e mantém observação até 300\,000 ms de estabilidade. O primeiro prazo dá tempo para defeitos de inicialização e operação imediata aparecerem \emph{antes} de remover a proteção do bootloader. O prazo total de 5 minutos estende a detecção para falhas que dependem de alguma atividade ou tempo de execução.

Após a confirmação de 2 minutos, o projeto usa ainda uma \emph{guarda} em memória RTC não inicializada (\texttt{RTC\_NOINIT\_ATTR}). Ela registra assinatura, endereço da partição em execução, contador de reinícios anormais e um valor de integridade calculado por XOR. Na partida, só aceita o registro se assinatura, integridade e endereço da partição coincidirem. Essa checagem evita interpretar resíduos de memória como estado válido ou atribuir a contagem a outro firmware. O estado não requer gravações frequentes na flash, preservando sua vida útil. A memória RTC costuma sobreviver a resets de software e de watchdog, mas \textbf{não é garantia diante de corte total de energia}; portanto, a contagem não deve ser apresentada como histórico persistente de falhas.

A guarda classifica como reinício anormal \texttt{ESP\_RST\_PANIC}, \texttt{ESP\_RST\_INT\_WDT}, \texttt{ESP\_RST\_TASK\_WDT} e \texttt{ESP\_RST\_WDT}. Após a confirmação inicial, cada partida por um desses motivos incrementa o contador. Ao atingir três, chama a API de invalidação e rollback. Reinícios intencionais ou de outra causa não entram nessa contagem. Ao completar o período estável de 5 minutos, a aplicação apaga o estado RTC da guarda. A guarda complementar não substitui o primeiro ensaio do bootloader, e o retorno depende de existir uma partição anterior válida. Se ela não existir, a API pode falhar; nesse caso, a recuperação exige diagnóstico e possível gravação por cabo.

\begin{nota}{Três significados de confirmação}
\textbf{Pedido aceito:} a interface HTTP ou o comando UART aceitou iniciar uma operação. \textbf{Imagem transferida:} os bytes foram gravados, o resumo e a versão conferem; no satélite, o mestre também viu a nova versão responder após o reboot. \textbf{Imagem confirmada como válida:} o aplicativo novo executou seus serviços essenciais por 2 minutos e chamou a confirmação do ESP-IDF. O acompanhamento de reinícios continua até 5 minutos. A palavra \emph{validada} em um log logo após a gravação refere-se à integridade do arquivo, não à aprovação funcional de 2 minutos.
\end{nota}

\section{Cancelamento, falhas e recuperação prática}
Antes de autorizar, \emph{Cancelar} descarta a oferta. Durante um download do mestre, a solicitação é verificada entre etapas de \texttt{esp\_https\_ota\_perform}; o gerenciador aborta antes de selecionar o slot novo. Durante transferência a satélite, o mestre para o pipeline e envia \emph{CANCELAR} com a sessão ativa; o satélite aborta o manipulador OTA e mantém o firmware atual. Um \emph{CANCELAR} de outra sessão é ignorado. Cancelar depois de o destino ter escolhido uma imagem e iniciado o reboot não é equivalente a desfazer uma instalação: nesse ponto, observar a versão que voltou a responder e o estado de rollback. A ação pode ser aceita antes de a tarefa encerrar; acompanhe o estado final e o erro em vez de repetir imediatamente comandos.

\begin{table}[htbp]\centering\small
\caption{Interpretação de falhas durante OTA}\label{tab:falhas-ota}
\begin{tabularx}{\textwidth}{@{}l X X@{}}\toprule
Indício & Causa provável a investigar & Primeira verificação\\\midrule
Manifesto ausente ou inválido & Release incompleto, URL inacessível, formato divergente ou Wi-Fi sem Internet. & Conferir sessão Wi-Fi, resposta HTTPS, \texttt{versao\_formato} e quatro anexos do release.\\
Alvo indisponível ou sem suporte & Telemetria antiga, UART desconectada, firmware satélite sem capacidade OTA. & Conferir GND comum, pinos RX/TX, versão e idade da última resposta.\\
Bloco repetido ou timeout & Perda de resposta UART, ruído, excesso de tráfego ou satélite ocupado gravando. & Conferir sequência, deslocamento, CRC, contadores de erro e cabos; não interpretar repetição como dupla gravação.\\
Imagem rejeitada ao finalizar & Tamanho, SHA-256, versão interna ou formato de aplicativo divergente. & Comparar binário e manifesto do mesmo release; recompilar o conjunto se necessário.\\
Versão nova não responde & Partida falhou, enlace não voltou ou rollback ocorreu. & Aguardar o prazo, ler logs seriais dos dois dispositivos e consultar versão real antes de repetir.\\
Volta à versão anterior & Falha no primeiro ensaio, serviço essencial ou três resets anormais na janela RTC. & Ler motivo de reset e erro de serviço; corrigir e publicar nova etiqueta.\\\bottomrule
\end{tabularx}\end{table}

Para instalar, comece com as três placas energizadas, GND comum e versões conhecidas. Confirme no painel da Equipe que mestre e satélites respondem. No terminal, abra OTA, verifique agora e anote versões instalada e oferecida. Autorize o alvo ou \emph{todos} dentro de 5 minutos. Observe bytes confirmados, taxa, erro e versão reportada após cada reinício; o terminal e a OLED do mestre apresentam o mesmo estado, incluindo uma barra de progresso. Aviso sonoro distingue atualização disponível, conclusão e falha. Aguarde a janela de 2 minutos para confirmação funcional e a de 5 minutos para encerramento da guarda antes de declarar a instalação estável. Se a oferta expirar, faça nova consulta; se um satélite estiver ausente, recupere o enlace antes de repetir.

O fluxo foi implementado com APIs OTA do ESP-IDF. Para detalhes do significado de \texttt{NEW}, \texttt{PENDING\_VERIFY}, \texttt{VALID}, \texttt{INVALID} e \texttt{ABORTED}, consulte a \href{https://docs.espressif.com/projects/esp-idf/en/latest/esp32s3/api-reference/system/ota.html}{documentação oficial de OTA do ESP-IDF}. Os procedimentos específicos do NEXUS, parâmetros de release e comandos estão registrados em \texttt{documentacao/ota.md}; os estados e prazos descritos aqui correspondem aos serviços e à guarda de rollback do repositório.
